\documentclass[a4paper]{article}
\usepackage{amsfonts}
\usepackage{amsmath}
\usepackage{amssymb}
\usepackage{graphicx}     

\begin{document}
  
\noindent \large Auteur: Stan Van Campenhout \\
\noindent \large Studierichting:  Informatica\\
\noindent \large Oefeningenreeks 6A nr. 25.1\\

\medskip

\normalsize

$ 1,2,5,10,17,26,37 $\\

Neem verschilrijen tot we aan een constate rij komen \\

$ 1,2,5,10,17,26,37 \Rightarrow 1,3,5,7,9,11 \Rightarrow 2, 2, 2, 2, 2 \Rightarrow $ orde 2 \\

$ x_n = an^2 + bn + c $\\

$ \Rightarrow c = 1, a+b+c = 2, 4a+2b+c = 5 \Rightarrow a + b = 1, 4a + 2b = 4 $\\ 

$ \Rightarrow a+b=1, 2a= 2\Rightarrow a = 1\Rightarrow b = 0 $\\

$ \Rightarrow x_n = n^2 + 1 $ met eerste term $ n = 0 $\\

$ \Rightarrow x_n = \left(n - 1\right)^2 + 1 = n^2-2n+2 $ met eerste term $ n = 1 $\\

$ \Rightarrow x_8 = 8^2-2\cdot8+2 = 64-16+2 = 50 $\\

\begin{thebibliography}{999}
%\bibitem{a} Referentie
\end{thebibliography}
\end{document} 
